\chapter{Metodologia e Desenvolvimento do Sistema}
\label{cap:metodologia}

Este capítulo detalha a metodologia de engenharia adotada na concepção, modelagem física e implementação do \textit{PID Simulator}. Apresentam-se os requisitos de projeto, a modelagem matemática das três plantas dinâmicas (veículo terrestre com relevo 3D, bola na barra com duplo integrador e tanque de nível de líquido com conservação volumétrica), a formulação da integração numérica em tempo real, as decisões arquiteturais do frontend em WebGL/Three.js com \textit{sandbox} segura de execução, a infraestrutura de serviços backend em FastAPI para emissão de relatórios e a conteinerização em Docker.

\section{Requisitos e Visão Geral da Arquitetura}
\label{sec:requisitos_arquitetura}

O projeto do sistema foi concebido para atender às demandas de um ambiente laboratorial virtual de ensino em Engenharia de Computação, aliando facilidade de uso a rigor numérico.

\subsection{Levantamento de Requisitos}
\label{subsec:levantamento_requisitos}

Os requisitos do sistema foram estruturados em categorias funcionais e não-funcionais:

\begin{itemize}
    \item \textbf{Requisitos Funcionais (RF):}
    \begin{itemize}
        \item \textbf{RF01 (Simulação em Malha Fechada):} Executar algoritmos de controle PID definidos programaticamente pelo usuário em tempo real sobre modelos físicos simulados;
        \item \textbf{RF02 (Editor de Código Integrado):} Prover ambiente de desenvolvimento em JavaScript com destaque sintático, auto-indentação e persistência local;
        \item \textbf{RF03 (Isolamento em Sandbox):} Executar o código do usuário em escopo protegido, interceptando exceções de execução sem interromper a interface gráfica;
        \item \textbf{RF04 (Bancada Tridimensional Interativa):} Renderizar a resposta cinemática e dinâmica da planta em cena 3D acelerada por hardware com câmera dinâmica;
        \item \textbf{RF05 (Osciloscópio Didático Multicanal):} Exibir a evolução temporal síncrona da variável de processo $y(t)$, referência $r(t)$, erro $e(t)$ e decomposição individual das forças $u_P(t), u_I(t), u_D(t)$;
        \item \textbf{RF06 (Assistente de Sintonia Automático):} Executar ensaio de degrau em malha aberta, calcular parâmetros FOPDT ($K, \tau, \theta$) e fornecer ganhos de Ziegler-Nichols e Cohen-Coon com aplicação direta no editor;
        \item \textbf{RF07 (Emissão de Relatórios Técnicos e Exportação):} Sintetizar relatórios em formato PDF com diagnóstico de controle e disponibilizar exportação de dados brutos em formato CSV para análise científica.
    \end{itemize}
    
    \item \textbf{Requisitos Não-Funcionais (RNF):}
    \begin{itemize}
        \item \textbf{RNF01 (Desempenho Temporal):} Manter taxa de atualização visual estável de no mínimo 30 quadros por segundo ($\text{FPS} \ge 30$) sob carga gráfica contínua;
        \item \textbf{RNF02 (Portabilidade e Acessibilidade):} Operar nativamente em navegadores modernos sem a necessidade de \textit{plugins} ou instalações locais proprietárias;
        \item \textbf{RNF03 (Determinismo Numérico):} Limitar o passo de integração temporal ($\Delta t \le 0{,}1\text{ s}$) para evitar divergência numérica em situações de oscilação severa;
        \item \textbf{RNF04 (Modularidade e Implantação):} Disponibilizar a aplicação encapsulada em contêineres Docker padronizados via `docker-compose`.
    \end{itemize}
\end{itemize}

\subsection{Arquitetura do Sistema}
\label{subsec:arquitetura_sistema}

A solução segue uma arquitetura cliente-servidor distribuída e desacoplada, esquematizada na Figura \ref{fig:arquitetura_geral}. O processamento de simulação numérica de alta frequência e a renderização gráfica ocorrem integralmente no cliente (\textit{browser-side}), aliviando o tráfego de rede e garantindo latência nula de feedback tátil. O servidor backend atua como repositório de persistência assíncrona, orquestrador de diagnósticos heurísticos e motor de compilação de relatórios técnicos em PDF.

\begin{figure}[htbp]
    \centering
    \caption{Arquitetura modular em camadas do PID Simulator integrando cliente WebGL e serviços backend.}
    \label{fig:arquitetura_geral}
    \fbox{\parbox{0.9\textwidth}{\centering
        \vspace{0.4cm}
        \textbf{[CAMADA CLIENTE --- NAVEGADOR WEB]} \\
        \vspace{0.2cm}
        $\boxed{\text{CodeMirror (Editor)}} \longleftrightarrow \boxed{\text{Sandbox JS (new Function)}} \longrightarrow \boxed{\text{PhysicsEngine (Euler explícito)}}$ \\
        \vspace{0.2cm}
        $\boxed{\text{PhysicsEngine}} \xrightarrow{\text{Estado Físico}} \boxed{\text{Scene3D (Three.js/WebGL)}} \quad \text{e} \quad \boxed{\text{ChartAnalytics (Chart.js)}}$ \\
        \vspace{0.4cm}
        $\Updownarrow \text{HTTP REST / JSON (fetch assíncrono)}$ \\
        \vspace{0.4cm}
        \textbf{[CAMADA SERVIDOR --- FASTAPI / PYTHON]} \\
        \vspace{0.2cm}
        $\boxed{\text{Endpoints REST (/runs)}} \longrightarrow \boxed{\text{ScoreService (Diagnóstico)}} \longrightarrow \boxed{\text{ReportGenerator (PDF/CSV)}}$ \\
        \vspace{0.2cm}
        $\boxed{\text{Endpoints REST}} \longleftrightarrow \boxed{\text{SQLAlchemy ORM}} \longleftrightarrow \boxed{\text{PostgreSQL Database}}$
        \vspace{0.4cm}
    }}
    \vspace{0.1cm}
    \fonte{Elaborada pelo autor (2026).}
\end{figure}

\section{Modelagem Física e Simulação Numérica}
\label{sec:modelagem_fisica}

Para expor o discente a desafios de controle de complexidade crescente, o sistema implementa três modelos físicos com características dinâmicas distintas.

\subsection{Planta 1: Dinâmica Longitudinal de Veículo com Relevo 3D}
\label{subsec:planta_carro}

A primeira planta modela a movimentação linear de um veículo terrestre de massa $m$ sujeito à força propulsiva do motor $u(t)$, atrito viscoso $b \cdot v(t)$ e perturbações gravitacionais induzidas pelo perfil topográfico da pista de testes. A equação diferencial de movimento de segundo grau em relação à posição $x(t)$ é dada por:

\begin{equation}
    m \, \frac{d^2 x(t)}{dt^2} + b \, \frac{dx(t)}{dt} = u(t) + F_{\text{relevo}}(x)
    \label{eq:dinamica_veiculo}
\end{equation}

A força gravitacional decorrente do relevo tridimensional é calculada dinamicamente a partir da inclinação instantânea $\alpha(x)$ da malha da pista:

\begin{equation}
    F_{\text{relevo}}(x) = - m \, g \, \sin(\alpha(x))
    \label{eq:forca_relevo}
\end{equation}

\noindent onde $\alpha(x) = \arctan\left(\frac{dh}{dx}\right)$ representa a inclinação tangencial derivada da função de elevação $h(x)$ do terreno no ponto ocupado pelo veículo. Esta formulação introduz uma perturbação não-linear de carga dependente da posição espacial, demandando ação integral ($u_I$) vigorosa para anular o erro de regime ao subir rampas.

\subsection{Planta 2: Dinâmica de Bola na Barra (Ball and Beam)}
\label{subsec:planta_ball_beam}

A segunda planta modela o clássico sistema de equilíbrio da bola na barra (\textit{Ball and Beam}), composto por uma esfera sólida de massa $m_b$, raio $R$ e momento de inércia $J_b = \frac{2}{5} m_b R^2$ que rola sem deslizamento sobre uma barra rígida de comprimento $L$ inclinada por um servomotor em torno de um pivô central. A equação dinâmica de movimento ao longo da barra em função do ângulo de inclinação $\alpha(t)$ é expressa por:

\begin{equation}
    \left( \frac{J_b}{R^2} + m_b \right) \ddot{x}(t) = m_b \, g \, \sin(\alpha(t)) - b_r \, \dot{x}(t) + d_{\text{pert}}(t)
    \label{eq:dinamica_ball_beam}
\end{equation}

Substituindo a inércia esférica e assumindo pequenas inclinações angulares ($\sin(\alpha) \approx \alpha$), obtém-se a aceleração linear da esfera:

\begin{equation}
    \ddot{x}(t) = \frac{5}{7} \, g \, \alpha(t) - \frac{b_r}{m_b} \, \dot{x}(t) + d_{\text{pert}}(t)
    \label{eq:aceleracao_ball_beam}
\end{equation}

\noindent na qual $x(t) \in [-L/2, +L/2]$ representa a posição da bola em relação ao centro e $\alpha(t)$ é a variável de atuação comandada pelo controlador. Sob a ótica de controle, o sistema configura um duplo integrador equivalente ($G(s) \approx \frac{K_{\text{bb}}}{s^2}$), sendo de natureza \textbf{intrinsecamente instável em malha aberta}. 

Essa característica impõe um desafio didático de grande relevância: controladores puramente proporcionais ($P$) ou proporcionais-integrais ($PI$) são incapazes de estabilizar a esfera, que oscila com amplitude crescente até atingir os batentes das extremidades. A estabilização demanda obrigatoriamente a ação derivativa ($K_d$), que atua antecipando o movimento e amortecendo a energia cinética da esfera.

\subsection{Planta 3: Dinâmica de Tanque de Nível de Líquido (Water Tank)}
\label{subsec:planta_tanque}

A terceira planta representa um processo contínuo de primeira ordem comum na indústria química, petrolífera e de saneamento: a regulação de nível em um reservatório cilíndrico de área transversal uniforme $A$, alimentado por uma bomba centrífuga com vazão de entrada $q_{\text{in}}(t)$ e escoamento gravitacional livre por um orifício de área $a_{\text{out}}$ na base. A dinâmica volumétrica é regida pelo princípio da conservação de massa:

\begin{equation}
    A \, \frac{dh(t)}{dt} = q_{\text{in}}(t) - q_{\text{out}}(t) - q_{\text{vaz}}(t)
    \label{eq:balanco_massa_tanque}
\end{equation}

Pela Lei de Torricelli, a vazão de escoamento no dreno é proporcional à raiz quadrada da coluna hidrostática:

\begin{equation}
    q_{\text{out}}(t) = a_{\text{out}} \sqrt{2 \, g \, h(t)}
    \label{eq:torricelli}
\end{equation}

Linearizando o termo de escoamento em torno de um nível de operação nominal $h_0$, obtém-se a resistência hidráulica equivalente $R_h = \frac{\sqrt{2 g h_0}}{g a_{\text{out}}}$ e a função de transferência clássica de primeira ordem:

\begin{equation}
    G(s) = \frac{H(s)}{Q_{\text{in}}(s)} = \frac{K_h}{\tau_h \, s + 1}
    \label{eq:ft_tanque}
\end{equation}

\noindent onde $\tau_h = A R_h$ é a constante de tempo do reservatório e $K_h = R_h$ é o ganho estático. O modelo implementa a não-linearidade física de saturação assimétrica do atuador (a bomba injeta líquido no intervalo $0 \le u \le 100\%$, porém é incapaz de aspirar líquido), permitindo ao aluno exercitar sintonias clássicas por Ziegler-Nichols e compensação de vazamentos de demanda $q_{\text{vaz}}(t)$.

\subsection{Discretização Numérica e Laço de Simulação}
\label{subsec:laco_simulacao}

A integração das equações diferenciais contínuas é realizada no motor de física (`PhysicsEngine.js`) através do método de Euler explícito acoplado ao ciclo de renderização do navegador via `requestAnimationFrame`. Em cada quadro $k$, calcula-se o intervalo temporal real decorrido $\Delta t = (t_k - t_{k-1})$:

\begin{align}
    a[k] &= \frac{u[k] + d[k] - b \, v[k-1]}{m} \label{eq:aceleracao_discreta} \\
    v[k] &= v[k-1] + a[k] \, \Delta t \label{eq:velocidade_discreta} \\
    x[k] &= x[k-1] + v[k] \, \Delta t \label{eq:posicao_discreta}
\end{align}

Para garantir estabilidade numérica mesmo sob eventuais atrasos de agendamento na \textit{thread} principal do navegador, o passo $\Delta t$ é rigidamente saturado em $\Delta t_{\max} = 0{,}1\text{ s}$.

\section{Implementação da Plataforma Frontend}
\label{sec:implementacao_frontend}

A interface de usuário foi concebida sob os princípios de ergonomia de estações de instrumentação, organizada em módulos complementares.

\subsection{Editor de Código e Sandbox de Execução Segura}
\label{subsec:editor_sandbox}

A programação do controlador é viabilizada através da integração do editor CodeMirror com suporte a realce sintático para JavaScript. O algoritmo definido pelo usuário é compilado em tempo de execução através do construtor `new Function`:

\begin{verbatim}
const userFunc = new Function(
    'error', 'integral', 'derivative', 'velocity', 'position', 'disturbance',
    customCode
);
const output = userFunc(error, integral, derivative, velocity, position, dist);
\end{verbatim}

Essa abordagem isola as variáveis locais do motor de física, impedindo a poluição do escopo global `window`. A chamada é interceptada por tratamento estrito de exceções; caso ocorra divisão por zero, retorno não-numérico (`NaN`) ou laço infinito acidental, a simulação é pausada preventivamente e uma mensagem descritiva é apresentada no console de diagnóstico da interface.

\subsection{Cena Tridimensional com Three.js}
\label{subsec:cena_threejs}

O ambiente tridimensional é estruturado sobre o WebGL utilizando a biblioteca Three.js. A pista de testes da Fase 1 é construída proceduralmente através de malhas segmentadas (`BufferGeometry`), permitindo a aplicação de deformações dinâmicas nos vértices de altura para gerar modos de relevo (plano, aclive de $+15^\circ$, declive de $-15^\circ$ e pista ondulada senoidal).

A câmera virtual acompanha a translação do veículo utilizando amortecimento exponencial (\textit{lerp}):

\begin{equation}
    \mathbf{P}_{\text{cam}}[k] = \mathbf{P}_{\text{cam}}[k-1] + \lambda \, (\mathbf{P}_{\text{alvo}}[k] - \mathbf{P}_{\text{cam}}[k-1])
    \label{eq:lerp_camera}
\end{equation}

\noindent onde $\lambda \in (0, 1)$ define o fator de suavização, garantindo uma transição visual contínua sem trepidação durante variações bruscas de velocidade.

\subsection{Osciloscópio Digital e Decomposição de Sinais}
\label{subsec:osciloscopio_chart}

O módulo de telemetria didática emprega a biblioteca Chart.js otimizada para atualização contínua em tempo real. O gráfico projeta simultaneamente seis curvas temporais correlacionadas:

\begin{enumerate}
    \item Variável de processo $y(t)$ (traço ciano);
    \item Sinal de referência $r(t)$ (linha tracejada vermelha);
    \item Parcela Proporcional $u_P(t) = K_p e(t)$ (traço âmbar pontilhado);
    \item Parcela Integral $u_I(t) = K_i \int e \, dt$ (traço esmeralda tracejado);
    \item Parcela Derivativa $u_D(t) = K_d \frac{de}{dt}$ (traço roxo traço-ponto);
    \item Resposta analítica teórica de referência $y_{\text{teorico}}(t)$ (linha cinza).
\end{enumerate}

Chaves seletoras (\textit{toggles}) na interface permitem ligar e desligar individualmente cada componente do esforço de controle, viabilizando ensaios analíticos onde o estudante observa o isolamento do efeito de amortecimento do termo $u_D$ ou a ação cumulativa do termo $u_I$.

\subsection{Módulo Automatizado de Sintonia Assistida}
\label{subsec:modulo_sintonia}

O componente `TuningAssistant.js` automatiza o procedimento experimental de sintonia em malha aberta:

\begin{enumerate}
    \item O usuário clica em ``Iniciar Ensaio de Malha Aberta'', disparando a aplicação de um degrau unitário constante $u(t) = 1{,}0$;
    \item O sistema monitora a resposta até a estabilização assintótica da velocidade ou posição;
    \item Algoritmicamente, o módulo localiza o ponto de maior derivada temporal $\left(\frac{dy}{dt}\right)_{\max}$, traça a reta tangente e determina a interseção com o eixo temporal, extraindo os parâmetros FOPDT: ganho estático $K$, tempo morto $\theta$ e constante de tempo $\tau$;
    \item As tabelas de Ziegler-Nichols e Cohen-Coon são computadas e apresentadas em uma janela modal com botão ``Aplicar Ganhos no Editor'', que injeta os valores calculados diretamente no código-fonte do estudante.
\end{enumerate}

\section{Infraestrutura Backend e Geração Científica}
\label{sec:backend_servicos}

A camada backend foi estruturada em Python utilizando o framework assíncrono FastAPI. A persistência de dados utiliza o SQLAlchemy ORM conectado a um banco relacional PostgreSQL, mapeando modelos para registros de ensaios laboratoriais (`RunModel`) e amostras temporais de telemetria.

\subsection{Diagnóstico Heurístico e Geração de Relatórios (PDF)}
\label{subsec:diagnostico_pdf}

Ao término de cada ensaio, a série temporal de telemetria é submetida ao endpoint `POST /api/v1/runs/`. O módulo de inteligência diagnóstica (`score_service.py`) analisa os dados consolidados e classifica a qualidade da resposta do controlador segundo regras heurísticas:

\begin{itemize}
    \item Se $M_p > 25\%$, o diagnóstico sinaliza ganho proporcional $K_p$ excessivo ou falta de ação derivativa $K_d$;
    \item Se $e_{ss} > 2\%$, recomenda-se a elevação do ganho integral $K_i$;
    \item Se o tempo de subida for excessivamente lento ($t_r > 3\tau$), indica-se sub-dimensionamento de $K_p$.
\end{itemize}

O gerador de documentos (`report_generator.py`), construído com as bibliotecas ReportLab e Matplotlib, sintetiza um laudo técnico completo em formato PDF contendo:
\begin{enumerate}[label=\alph*)]
    \item Cabeçalho institucional com carimbo de data, hora e fase experimental;
    \item Tabela de métricas consolidadas ($t_r, t_s, M_p, e_{ss}, \text{RMSE}, \text{IAE}, \text{ITAE}$);
    \item Painel superior com curvas comparativas de saída $y(t)$, referência $r(t)$ e erro $e(t)$;
    \item Painel inferior com a decomposição vetorial de esforços $u_P(t), u_I(t), u_D(t)$ e limites de saturação;
    \item Parecer técnico detalhado com recomendações de ajuste paramétrico.
\end{enumerate}

Adicionalmente, o endpoint `GET /api/v1/runs/{id}/export-csv` gera conjuntos de dados estruturados com as séries temporais brutas, permitindo que os discentes importem os resultados em ferramentas científicas (como Python SciPy, MATLAB ou OriginLab) para confecção de artigos e relatórios de laboratório.

\section{Conteinerização e Implantação com Docker}
\label{sec:conteinerizacao}

Para assegurar reprodutibilidade experimental estrita e eliminar conflitos de dependências em estações de trabalho de laboratório, toda a aplicação foi empacotada em contêineres Docker orquestrados por `docker-compose`. 

O serviço frontend opera sob servidor web Nginx ou servidor estático leve, expondo a porta `8080`, enquanto a API backend FastAPI executa via Uvicorn na porta `8081` com migrações automáticas de banco de dados. Um único comando (`docker compose up --build`) é suficiente para instanciar integralmente o laboratório virtual em qualquer sistema operacional (Linux, macOS, Windows).
